\chapter{Lesson 6: Self-Introduction (\ja{XはYです}, \ja{の}, \ja{か})}

\section{Vocabulary}

\subsection{People and school}
\begin{center}
\begin{tabular}{lll}
\toprule
Japanese & Romaji & Meaning \\
\midrule
\ja{わたし} & watashi & I \\
\ja{あなた} & anata & you (avoid; use the person's name) \\
\ja{がくせい} & gakusei & student \\
\ja{だいがくせい} & daigakusei & university student \\
\ja{りゅうがくせい} & ryūgakusei & international student \\
\ja{せんせい} & sensei & teacher, professor \\
\ja{だいがく} & daigaku & university \\
\ja{せんこう} & senkō & major \\
\ja{〜ねんせい} & \ldots nensei & \ldots-year student \\
\ja{ともだち} & tomodachi & friend \\
\ja{なまえ} & namae & name \\
\ja{でんわばんごう} & denwa bangō & telephone number \\
\ja{〜じん} & \ldots jin & person from \ldots \\
\ja{〜ご} & \ldots go & \ldots\ language \\
\ja{はい} / \ja{ええ} & hai / ee & yes \\
\ja{いいえ} & iie & no \\
\bottomrule
\end{tabular}
\end{center}

\subsection{Countries, nationalities and languages}
\begin{center}
\begin{tabular}{llll}
\toprule
Country & Person & Language & Meaning \\
\midrule
\ja{にほん} Nihon & \ja{にほんじん} & \ja{にほんご} & Japan \\
\ja{カナダ} Kanada & \ja{カナダじん} & --- & Canada \\
\ja{アメリカ} Amerika & \ja{アメリカじん} & \ja{えいご} (English) & USA \\
\ja{ちゅうごく} Chūgoku & \ja{ちゅうごくじん} & \ja{ちゅうごくご} & China \\
\ja{かんこく} Kankoku & \ja{かんこくじん} & \ja{かんこくご} & Korea \\
\ja{インド} Indo & \ja{インドじん} & \ja{ヒンディーご} & India \\
\ja{イギリス} Igirisu & \ja{イギリスじん} & \ja{えいご} & UK \\
\ja{フランス} Furansu & \ja{フランスじん} & \ja{フランスご} & France \\
\bottomrule
\end{tabular}
\end{center}

\subsection{Majors (\ja{せんこう})}
\begin{center}
\begin{tabular}{lll@{\qquad}lll}
\toprule
\ja{すうがく} & sūgaku & math & \ja{けいざい} & keizai & economics \\
\ja{コンピューター} & konpyūtā & computer science & \ja{ぶんがく} & bungaku & literature \\
\ja{こうがく} & kōgaku & engineering & \ja{れきし} & rekishi & history \\
\ja{ビジネス} & bijinesu & business & \ja{かがく} & kagaku & science \\
\ja{いがく} & igaku & medicine & \ja{せいぶつがく} & seibutsugaku & biology \\
\bottomrule
\end{tabular}
\end{center}

\section{Grammar 1: \ja{XはYです} ``X is Y''}

\begin{center}
\begin{tikzpicture}[b/.style={draw, rounded corners, minimum height=0.9cm, minimum width=2.4cm},
    p/.style={draw, circle, fill=orange!15, minimum size=0.9cm}]
  \node[b, fill=blue!6] (x) at (0,0) {X (topic)};
  \node[p] (wa) at (2.1,0) {\ja{は}};
  \node[b, fill=green!8] (y) at (4.2,0) {Y (noun)};
  \node[b, fill=gray!10, minimum width=1.4cm] (d) at (6.4,0) {\ja{です}};
  \node[font=\small] at (0,-0.85) {\ja{わたし}};
  \node[font=\small] at (2.1,-0.85) {wa};
  \node[font=\small] at (4.2,-0.85) {\ja{がくせい}};
  \node[font=\small] at (6.4,-0.85) {desu};
  \node[font=\small, text width=3.2cm, align=center] at (0,1.15) {``As for X\ldots''};
  \node[font=\small, text width=3.2cm, align=center] at (5.3,1.15) {``\ldots(it) is Y.''};
\end{tikzpicture}
\end{center}

\begin{itemize}
  \item \ja{は} (read \emph{wa}) marks the \emph{topic}: what the sentence is about.
  \item \ja{です} is the polite copula ``is/am/are''. It does not change for person or
        number.
  \item Japanese has no articles (a/the) and nouns do not mark plural.
\end{itemize}

\begin{example}
\leavevmode
\begin{itemize}
  \item \ja{わたしはりゅうがくせいです。} \emph{Watashi wa ryūgakusei desu.}
        I am an international student.
  \item \ja{やまださんはせんせいです。} \emph{Yamada-san wa sensei desu.}
        Ms./Mr.\ Yamada is a teacher.
  \item \ja{せんこうはすうがくです。} \emph{Senkō wa sūgaku desu.} (My) major is math.
  \item \ja{いま、にじです。} \emph{Ima, niji desu.} It's two o'clock now.
\end{itemize}
\end{example}

\begin{important}[Drop what is understood]
If the topic is clear from context, leave it out. Saying \ja{わたしは} in every sentence
sounds unnatural and self-centred. After the first sentence of a self-introduction,
just say \ja{いちねんせいです。} ``(I'm) a first-year student.''
\end{important}

\subsection{Negative: \ja{じゃないです}}
\begin{center}
\begin{tabular}{lll}
\toprule
 & Form & Example \\
\midrule
affirmative & Y \ja{です} & \ja{がくせいです} \\
negative (spoken) & Y \ja{じゃないです} & \ja{がくせいじゃないです} \\
negative (formal/written) & Y \ja{ではありません} & \ja{がくせいではありません} \\
\bottomrule
\end{tabular}
\end{center}
\begin{example}
\ja{わたしはにほんじんじゃないです。カナダじんです。}
\emph{Watashi wa Nihonjin ja nai desu. Kanadajin desu.} I'm not Japanese. I'm Canadian.
\end{example}

\section{Grammar 2: questions with \ja{か}}

Add \ja{か} to the end of any statement to make a yes/no question. Word order does not
change. In writing, \ja{。} (not ?) is the standard punctuation.
\begin{example}
\leavevmode
\begin{itemize}
  \item[Q:] \ja{リーさんはちゅうごくじんですか。} \emph{Rī-san wa Chūgokujin desu ka.}
        Are you (Mr.\ Lee) Chinese?
  \item[A:] \ja{はい、そうです。} \emph{Hai, sō desu.} Yes, that's right.
  \item[A:] \ja{いいえ、かんこくじんです。} \emph{Iie, Kankokujin desu.} No, I'm Korean.
\end{itemize}
\end{example}

\subsection{Question words}
For information questions, put a question word where the answer would go and keep
\ja{か} at the end.
\begin{center}
\begin{tabular}{lll}
\toprule
Question word & Meaning & Example \\
\midrule
\ja{なん} / \ja{なに} & what & \ja{せんこうはなんですか。} What is your major? \\
\ja{だれ} & who & \ja{あのひとはだれですか。} Who is that person? \\
\ja{どこ} & where & \ja{トイレはどこですか。} Where is the toilet? \\
\ja{いつ} & when & \ja{しけんはいつですか。} When is the exam? \\
\ja{なんさい} & how old & \ja{なんさいですか。} How old are you? \\
\ja{なんねんせい} & what year & \ja{なんねんせいですか。} What year are you in? \\
\ja{なんじ} & what time & \ja{いまなんじですか。} What time is it now? \\
\bottomrule
\end{tabular}
\end{center}

\begin{remark}
\ja{なに} becomes \ja{なん} before \ja{です}, before counters (\ja{なんじ},
\ja{なんさい}) and before d-, n-, t-sounds. \ja{なにご} ``what language'' keeps
\ja{なに}.
\end{remark}

\section{Grammar 3: the particle \ja{の}}

\ja{AのB} links two nouns; A describes or limits B. Most often it shows possession or
affiliation, but the relationship can be many things. The \emph{main} noun is the last
one (B).

\begin{center}
\begin{tikzpicture}[b/.style={draw, rounded corners, minimum height=0.8cm, minimum width=2.2cm}]
  \node[b] (a) at (0,0) {\ja{ウォータールーだいがく}};
  \node[circle, draw, fill=orange!15] (no) at (2.6,0) {\ja{の}};
  \node[b, fill=green!8, minimum width=1.8cm] (b) at (4.4,0) {\ja{がくせい}};
  \draw[-Stealth, thick] (a.south) to[bend right=30] node[below, font=\small] {modifies} (b.south);
  \node[font=\small] at (0,0.75) {A: Waterloo University};
  \node[font=\small] at (4.4,0.75) {B: student (head)};
\end{tikzpicture}
\end{center}

\begin{example}[Uses of \ja{の}]
\leavevmode
\begin{itemize}
  \item \ja{わたしのなまえ} \emph{watashi no namae} my name (possession)
  \item \ja{だいがくのせんせい} \emph{daigaku no sensei} a university professor (affiliation)
  \item \ja{にほんごのほん} \emph{Nihongo no hon} a Japanese-language book (content)
  \item \ja{にほんのくるま} \emph{Nihon no kuruma} a Japanese car (origin)
  \item \ja{わたしのともだちのでんわばんごう} my friend's phone number (\ja{の} can chain)
\end{itemize}
\end{example}

\ja{の} can also stand in for a noun already understood:
\ja{これはだれのですか。} ``Whose is this?'' --- \ja{やまださんのです。}
``It's Yamada-san's.''

\section{Putting it together: a self-introduction (\ja{じこしょうかい})}

A standard self-introduction opens with \ja{はじめまして}, gives name, origin, school
year and major, and closes with \ja{よろしくおねがいします}, with a bow.

\begin{example}[My self-introduction from the oral quiz]
\leavevmode
\begin{center}
\begin{tabular}{ll}
\ja{はじめまして。} & Nice to meet you. \\
\ja{グプタです。} & I'm Gupta. \\
\ja{カナダじんです。} & I'm Canadian. \\
\ja{ウォータールーだいがくのいちねんせいです。} & I'm a first-year at UW. \\
\ja{せんこうはコンピューターです。} & My major is computer science. \\
\ja{じゅうはっさいです。} & I'm 18. \\
\ja{どうぞよろしくおねがいします。} & Pleased to meet you. \\
\end{tabular}
\end{center}
\end{example}

\begin{aside}
\textbf{Introducing yourself in Japan.} Give your family name (or full name, family
first if you want to be thorough). In business settings this is when business cards
(\ja{めいし}, \emph{meishi}) are exchanged: offered and received with both hands, read
carefully, and placed on the table in front of you, never shoved in a pocket. For
foreign names, it is fine to keep your own name order in katakana, but some people
reverse it to match Japanese convention; either way, say which part is the family name.
\end{aside}

\section*{Lesson 6 review}
\begin{itemize}
  \item \ja{XはYです}: X is Y. Negative \ja{じゃないです}/\ja{ではありません}.
  \item \ja{か} at the end makes a question; answer \ja{はい、そうです} or
        \ja{いいえ、〜じゃないです}.
  \item Question words go where the answer goes: \ja{なん}, \ja{だれ}, \ja{どこ},
        \ja{いつ}.
  \item \ja{AのB}: B described by A; the last noun is the head.
  \item Drop \ja{わたしは} once the topic is clear.
\end{itemize}
